\chapter{Backend Architecture, HTTP Evolution, and API Paradigms}
\label{chap:backend-arch-http-api}

\begin{chaptergoalsbox}
\textbf{What You Will Learn:}
\begin{itemize}[leftmargin=12pt,itemsep=2pt]
  \item The architectural evolution of HTTP: HTTP/1.1 pipelining limitations, HTTP/2 binary framing and HPACK, and HTTP/3 QUIC stream independence over UDP.
  \item Roy Fielding's REST constraints and the Richardson Maturity Model up to HATEOAS level 3.
  \item The WebSocket protocol (RFC 6455): 101 Switching Protocols upgrade handshakes, client XOR frame masking, and horizontal state scaling.
  \item Server-Sent Events (SSE): unidirectional \texttt{text/event-stream} mechanics, browser recovery via \texttt{Last-Event-ID}, and HTTP/2 transport efficiency.
  \item gRPC and Protocol Buffers v3: binary varints, zigzag encoding, field wire types, and HTTP/2 trailer-based status framing.
\end{itemize}

\textbf{What You Should Already Understand:}
\begin{itemize}[leftmargin=12pt,itemsep=2pt]
  \item The TCP/IP 4-layer model and three-way handshake (\texttt{SYN}, \texttt{SYN-ACK}, \texttt{ACK}).
  \item Basic HTTP request and response header semantics.
\end{itemize}

\textbf{Why This Topic Exists:}
Distributed backend engineering is fundamentally the study of network communication between decoupled services. Choosing the wrong transport protocol, serialization format, or API paradigm degrades throughput, inflates tail latency, and creates cascading failure modes across microservice fleets. Senior engineers must understand protocol mechanics at the wire level to design resilient communication fabrics.

\textbf{Motivating Practical Problem:}
A financial market data trading portal streams continuous stock quotes and order book updates to 200,000 active web and mobile clients while exchanging trade execution commands with internal clearing microservices. Naive polling over HTTP/1.1 collapses origin load balancers with connection overhead, while attempting to run internal RPC over REST JSON causes excessive CPU serialization overhead. To resolve this, the architecture must partition workloads across HTTP/2 multiplexing, low-overhead WebSockets, Server-Sent Events, and sub-millisecond gRPC Protobuf pipelines.
\end{chaptergoalsbox}

\section{The Evolution of HTTP: From HTTP/1.1 to HTTP/3 QUIC}

The Hypertext Transfer Protocol (HTTP) has undergone three decades of rigorous re-engineering to keep pace with modern network scale.

\subsection{HTTP/1.1: Foundations and Application Head-of-Line Blocking}

Standardized in RFC 2616 (and revised in RFC 7230--7235), HTTP/1.1 introduced persistent connections via the \texttt{Connection: keep-alive} header, enabling reuse of an underlying TCP connection for multiple sequential requests.

\subsubsection{The Flaw of HTTP/1.1 Pipelining and Head-of-Line Blocking}
To reduce latency, HTTP/1.1 introduced pipelining, allowing clients to send multiple requests without waiting for intermediate responses. However, HTTP/1.1 requires the server to **send responses in the exact order requests were received**. If request 1 requires an expensive 3-second database query, requests 2, 3, and 4 are blocked behind it on that TCP socket. This phenomenon is known as **Application-Layer Head-of-Line (HoL) Blocking**. 

Because packet loss and proxy implementations frequently corrupted pipelined responses, modern browsers universally disabled pipelining, opening 6 parallel TCP connections per host instead—wasting server memory and kernel socket buffers.

\subsection{HTTP/2: Binary Framing, Streams, and HPACK Compression}

Standardized in RFC 7540, HTTP/2 revolutionized transport by introducing a **binary framing layer** between the socket and the application.

\begin{figure}[htbp]
\centering
\begin{tikzpicture}[
  scale=0.9, transform shape,
  node distance=0.8cm and 1.2cm,
  layer/.style={rectangle, draw=headingblue, fill=codebg, rounded corners=2pt, thick, align=center, inner sep=6pt, font=\small\sffamily},
  accent/.style={rectangle, draw=intuitionTeal, fill=intuitionBg, rounded corners=2pt, thick, align=center, inner sep=6pt, font=\small\sffamily\bfseries},
  line/.style={draw=headingblue, -{Stealth[length=2.5mm]}, thick}
]
  \node[layer, minimum width=10cm] (app) {Application Layer (HTTP Methods, Headers, Status Codes)};
  \node[accent, below=of app, minimum width=10cm] (frame) {HTTP/2 Binary Framing Layer\\(Multiplexed Streams, HEADERS \& DATA Frames)};
  \node[layer, below=of frame, minimum width=10cm] (tcp) {Transport Layer: Single Persistent TCP Connection};

  \draw[line] (app) -- (frame);
  \draw[line] (frame) -- (tcp);
\end{tikzpicture}
\caption{HTTP/2 binary framing layer enabling stream multiplexing over a single TCP connection.}
\label{fig:http2-framing}
\end{figure}

\subsubsection{Streams and Multiplexing}
HTTP/2 decomposes requests and responses into discrete binary frames:
\begin{itemize}[leftmargin=15pt]
  \item \textbf{Stream:} An independent, bi-directional sequence of frames exchanged across the connection. Each stream has an integer identifier (client-initiated streams use odd numbers; server-initiated push streams use even numbers).
  \item \textbf{Frame Types:} Fixed 9-byte header followed by variable payload. Types include \texttt{HEADERS} (opens a stream with compressed headers), \texttt{DATA} (carries body bytes), \texttt{SETTINGS}, \texttt{RST\_STREAM} (cancels a stream without closing the connection), and \texttt{WINDOW\_UPDATE} (per-stream and connection flow control).
\end{itemize}

Frames from dozens of concurrent requests are interleaved arbitrarily across a single TCP connection, eliminating application-layer Head-of-Line blocking completely!

\subsubsection{HPACK Header Compression (RFC 7541)}
In HTTP/1.1, headers (cookies, user-agents, tokens) are transmitted as uncompressed ASCII strings on every request, wasting kilobytes of bandwidth. HTTP/2 introduced **HPACK**:
\begin{enumerate}[leftmargin=15pt]
  \item \textbf{Static Table:} A pre-defined table of 61 common headers (e.g., index 2 maps to \texttt{:method: GET}, index 7 maps to \texttt{:scheme: https}). Common headers are transmitted as single-byte integers.
  \item \textbf{Dynamic Table:} State maintained across the connection. New headers (e.g., \texttt{Authorization: Bearer <jwt>}) are emitted once, stored in the dynamic table, and referenced by index in subsequent frames.
  \item \textbf{Huffman Coding:} Static Huffman tree optimizes custom strings.
\end{enumerate}

\subsubsection{The Fatal Limitation: TCP-Level Head-of-Line Blocking}
Although HTTP/2 solves application HoL blocking, it remains bound to a single TCP stream. TCP guarantees strictly ordered byte streams. If a single network packet is dropped on the physical wire, the operating system kernel holds all subsequent packets in its TCP receive buffer until the dropped packet is retransmitted (via TCP SACK/ACK). Under high packet loss (e.g., 2\% packet loss on cellular networks), **every multiplexed HTTP/2 stream on that connection freezes**, performing worse than multiple HTTP/1.1 TCP connections!

\subsection{HTTP/3: QUIC over UDP and Zero-RTT Reconnection}

Standardized in RFC 9114, HTTP/3 replaces TCP entirely with **QUIC (RFC 9000)**, an encrypted, multiplexed transport protocol built on top of **UDP**.

\begin{figure}[htbp]
\centering
\begin{tikzpicture}[
  scale=0.9, transform shape,
  node distance=0.8cm and 1.2cm,
  proto/.style={rectangle, draw=headingblue, fill=codebg, rounded corners=2pt, thick, align=center, inner sep=6pt, font=\small\sffamily},
  accent/.style={rectangle, draw=warningRed, fill=warningBg, rounded corners=2pt, thick, align=center, inner sep=6pt, font=\small\sffamily\bfseries},
  line/.style={draw=headingblue, -{Stealth[length=2.5mm]}, thick}
]
  % Stack HTTP/2
  \node[proto, minimum width=5cm] (h2) {\textbf{HTTP/2}};
  \node[proto, below=0.3cm of h2, minimum width=5cm] (tls12) {TLS 1.2 / 1.3};
  \node[proto, below=0.3cm of tls12, minimum width=5cm] (tcp) {TCP (Kernel)};

  % Stack HTTP/3
  \node[accent, right=2.0cm of h2, minimum width=5cm] (h3) {\textbf{HTTP/3 (QPACK)}};
  \node[accent, below=0.3cm of h3, minimum width=5cm] (quic) {QUIC (Streams + Congestion + TLS 1.3)};
  \node[accent, below=0.3cm of quic, minimum width=5cm] (udp) {UDP (Kernel / Userland)};

  \node[above=0.2cm of h2, font=\small\sffamily\bfseries] {HTTP/2 Protocol Stack};
  \node[above=0.2cm of h3, font=\small\sffamily\bfseries] {HTTP/3 Protocol Stack};
\end{tikzpicture}
\caption{Protocol stack comparison: HTTP/2 over TCP+TLS vs HTTP/3 over QUIC+UDP.}
\label{fig:http3-stack}
\end{figure}

\subsubsection{Key Capabilities of HTTP/3 QUIC}
\begin{enumerate}[leftmargin=15pt]
  \item \textbf{Independent Stream Transmission:} QUIC implements streams natively in userspace over UDP datagrams. If packet loss affects Stream 4, Streams 1, 2, 3, and 5 continue processing without a single millisecond of delay. TCP Head-of-Line blocking is eradicated.
  \item \textbf{Integrated TLS 1.3 Handshake:} In HTTP/2 over TCP, connection establishment requires a TCP 3-way handshake (1 RTT) followed by a TLS 1.3 handshake (1 RTT), totaling 2 RTTs before data flows. QUIC fuses the transport and cryptographic handshake into a **single 1-RTT exchange** (or **0-RTT** on repeated sessions).
  \item \textbf{Connection Migration:} TCP connections are bound to a 4-tuple: \texttt{(src\_ip, src\_port, dst\_ip, dst\_port)}. When a mobile user switches from Wi-Fi to cellular LTE, their IP changes, immediately breaking all open TCP connections. QUIC identifies connections via a random 64-bit **Connection ID (CID)** independent of IP or port. Sockets transition seamlessly without reconnecting!
\end{enumerate}

% ==========================================
\section{REST Architectural Constraints and the Richardson Maturity Model}

Introduced in Roy Fielding's 2000 doctoral dissertation, **Representational State Transfer (REST)** is not a protocol or standard, but an architectural style governed by six strict constraints:

\begin{enumerate}[leftmargin=15pt]
  \item \textbf{Client-Server Separation:} Separation of concerns; UI/client concerns are decoupled from data persistence.
  \item \textbf{Statelessness:} Each request from client to server must contain all of the information necessary to understand and complete the request. No session state is held on the server between requests.
  \item \textbf{Cacheability:} Responses must explicitly define whether they are cacheable (\texttt{Cache-Control}, \texttt{ETag}).
  \item \textbf{Uniform Interface:} The bedrock of REST, comprising:
  \begin{itemize}
    \item Resource identification in requests (URIs).
    \item Manipulation of resources through representations (JSON, XML).
    \item Self-descriptive messages (MIME types, standard status codes).
    \item Hypermedia as the Engine of Application State (HATEOAS).
  \end{itemize}
  \item \textbf{Layered System:} Clients cannot tell whether they are connected directly to the end server or to an intermediary proxy, CDN, or load balancer.
  \item \textbf{Code on Demand (Optional):} Servers can temporarily extend client functionality by transferring executable code (e.g., JavaScript).
\end{enumerate}

\subsection{The Richardson Maturity Model}

Leonard Richardson formalized an analytical model evaluating how closely a web API adheres to RESTful principles:

\begin{table}[htbp]
\centering
\small\sffamily
\begin{tabular}{@{}lll@{}}
\toprule
\textbf{Level} & \textbf{Core Characteristic} & \textbf{Typical Implementation Pattern} \\ \midrule
\textbf{Level 0} & The Swamp of POX & Single URI (e.g., \texttt{/api}), single method (\texttt{POST}), XML/JSON RPC. \\
\textbf{Level 1} & Resources & Discrete URIs per entity (\texttt{/orders/123}), but still using arbitrary methods. \\
\textbf{Level 2} & HTTP Verbs & Standard HTTP verbs (\texttt{GET}, \texttt{POST}, \texttt{PUT}, \texttt{DELETE}) and status codes (201, 404). \\
\textbf{Level 3} & HATEOAS & Hypermedia controls embedded in representations driving client state transitions. \\ \bottomrule
\end{tabular}
\caption{The Richardson Maturity Model for RESTful APIs.}
\label{tab:richardson}
\end{table}

\subsection{Safe vs. Idempotent Methods}

\begin{itemize}[leftmargin=15pt]
  \item \textbf{Safe Methods (\texttt{GET}, \texttt{HEAD}, \texttt{OPTIONS}):} Must not alter server state. Invoking a safe method 1 time or 10,000 times leaves the system state unchanged. Caches may freely pre-fetch and store safe responses.
  \item \textbf{Idempotent Methods (\texttt{PUT}, \texttt{DELETE}, \texttt{GET}):} Executing the operation $N$ times produces the exact same server state as executing it 1 time:
  $$f(f(x)) = f(x)$$
  Replacing a resource via \texttt{PUT /users/42} with body \texttt{\{"name": "Alice"\}} produces the same state regardless of network retries.
  \item \textbf{Non-Idempotent Methods (\texttt{POST}, \texttt{PATCH}):} Submitting \texttt{POST /orders} repeatedly creates multiple orders and charges the credit card multiple times unless protected by application idempotency keys.
\end{itemize}

% ==========================================
\section{WebSockets (RFC 6455): Full-Duplex Bi-Directional Streaming}

When applications require sub-millisecond, bi-directional event delivery—such as online trading portals, collaborative document editing, or multiplayer gaming—standard HTTP request-response overhead becomes prohibitive. **WebSockets (RFC 6455)** provides full-duplex communication over a single persistent TCP socket.

\subsection{The HTTP 101 Switching Protocols Handshake}

A WebSocket connection begins as a standard HTTP/1.1 request:

\begin{lstlisting}[language=sh,caption={Client Handshake Request}]
GET /chat HTTP/1.1
Host: server.example.com
Upgrade: websocket
Connection: Upgrade
Sec-WebSocket-Key: dGhlIHNhbXBsZSBub25jZQ==
Sec-WebSocket-Version: 13
\end{lstlisting}

The server validates the request and proves it understands RFC 6455 by computing:
$$\text{AcceptKey} = \text{Base64}(\text{SHA-1}(\text{Sec-WebSocket-Key} \; \Vert \; \text{"258EAFA5-E914-47DA-95CA-C5AB0DC85B11"}))$$
The constant GUID is hardcoded in RFC 6455. The server responds with HTTP status 101:

\begin{lstlisting}[language=sh,caption={Server Handshake Response}]
HTTP/1.1 101 Switching Protocols
Upgrade: websocket
Connection: Upgrade
Sec-WebSocket-Accept: s3pPLMBiTxaQ9kYGzzhZRbK+xOo=
\end{lstlisting}
Once the 101 response is transmitted, the HTTP parser is detached, and both sides transition to framing binary WebSocket packets.

\subsection{WebSocket Frame Structure and Client Masking}

\begin{figure}[htbp]
\centering
\begin{tikzpicture}[
  scale=0.9, transform shape,
  node distance=0.1cm,
  fld/.style={rectangle, draw=headingblue, fill=codebg, thick, minimum height=0.9cm, align=center, font=\scriptsize\ttfamily}
]
  \node[fld, minimum width=1.0cm] (fin) {FIN\\(1b)};
  \node[fld, minimum width=1.0cm, right=of fin] (rsv) {RSV\\(3b)};
  \node[fld, minimum width=1.5cm, right=of rsv] (op) {Opcode\\(4b)};
  \node[fld, minimum width=1.2cm, right=of op] (mask) {MASK\\(1b)};
  \node[fld, minimum width=2.2cm, right=of mask] (len) {Payload Len\\(7 bits)};
  \node[fld, minimum width=3.0cm, right=of len] (mkey) {Masking Key\\(0 or 4 bytes)};
  \node[fld, minimum width=3.5cm, below=0.3cm of fin, xshift=5.0cm] (payload) {Extension Data and Application Payload};

  \node[draw=intuitionTeal, dashed, fit=(fin) (rsv) (op) (mask) (len) (mkey) (payload), inner sep=4pt] {};
\end{tikzpicture}
\caption{The WebSocket binary frame structure (RFC 6455).}
\label{fig:websocket-frame}
\end{figure}

\subsubsection{Why Client-to-Server Masking is Mandatory}
RFC 6455 strictly mandates that **every frame sent from client to server MUST be masked using a 4-byte random XOR key**. If an unmasked frame is received by a compliant server, the server must terminate the connection immediately.

\begin{underthehoodbox}[The Cache Poisoning Exploit That Mandated Masking]
Before mandatory masking, malicious JavaScript in a browser could send an unmasked WebSocket frame containing raw HTTP headers (\texttt{GET /evil.js HTTP/1.1...}). Vulnerable transparent proxies intercepted this frame, mistook it for an authentic HTTP request from the victim, and cached malicious attacker payloads for innocent users. 

XOR masking scrambles bytes using a dynamic 4-byte key generated per frame. Intermediary proxies see only random white noise, rendering proxy cache poisoning attacks impossible.
\end{underthehoodbox}

\subsubsection{Unmasking Payload Algorithm}
For byte index $i$ in payload of length $N$:
$$P_{\text{unmasked}}[i] = P_{\text{masked}}[i] \oplus K[i \pmod 4]$$

% ==========================================
\section{Server-Sent Events (SSE): Unidirectional Streaming}

When communication flows exclusively from server to client (e.g., real-time notifications, live financial tickers, LLM streaming token responses), WebSockets introduce unnecessary complexity. **Server-Sent Events (SSE)** provides a lightweight, unidirectional transport built directly on HTTP.

\subsection{The \texttt{text/event-stream} Protocol}

An SSE connection is initiated as a standard HTTP GET request with header \texttt{Accept: text/event-stream}. The server keeps the HTTP connection open indefinitely, emitting formatted UTF-8 text records:

\begin{lstlisting}[language=sh,caption={SSE Wire Format Stream}]
HTTP/1.1 200 OK
Content-Type: text/event-stream
Cache-Control: no-cache
Connection: keep-alive

event: stock-price
id: 10452
data: {"symbol": "AAPL", "price": 182.45}

event: stock-price
id: 10453
data: {"symbol": "GOOGL", "price": 141.20}

\end{lstlisting}

Every event record terminates with two consecutive newline characters (\texttt{\textbackslash n\textbackslash n}).

\subsection{Automatic Browser Reconnection via \texttt{Last-Event-ID}}

If the network connection drops:
\begin{enumerate}[leftmargin=15pt]
  \item The browser's native \texttt{EventSource} client automatically reconnects after a configurable backoff (set via the \texttt{retry: 5000} field).
  \item The browser includes the HTTP header \texttt{Last-Event-ID: 10453}, identifying the last successfully processed event.
  \item The backend reads this header and replays any missing messages from a message broker or Redis stream, achieving seamless zero-data-loss recovery!
\end{enumerate}

\begin{intuitionbox}[SSE vs. WebSockets: When to Choose Which?]
\begin{itemize}[leftmargin=15pt]
  \item \textbf{Use Server-Sent Events (SSE) when:} Traffic is strictly unidirectional (server $\rightarrow$ client), such as LLM text completion streams, news feeds, and live dashboard metrics. SSE runs natively over HTTP/2, sharing the single multiplexed TCP/QUIC socket without firewall blocking.
  \item \textbf{Use WebSockets when:} Traffic is truly bidirectional (client $\leftrightarrow$ server) with high frequency, such as real-time collaborative whiteboards, chat apps, and competitive multiplayer games.
\end{itemize}
\end{intuitionbox}

% ==========================================
\section{gRPC and Protocol Buffers: Microservice Wire Formats}

In high-throughput microservice architectures, REST over JSON introduces severe performance penalties:
\begin{itemize}[leftmargin=15pt]
  \item JSON is human-readable text requiring computationally expensive string parsing, floating-point parsing, and memory allocations.
  \item Field keys (e.g., \texttt{"transaction\_identification\_uuid"}) are repeated on every single message across the wire.
\end{itemize}
**gRPC** (Google Remote Procedure Call) solves this using **Protocol Buffers (Protobuf)** serialization over HTTP/2 transport.

\subsection{Protobuf v3 Wire Format Mechanics}

Protobuf serializes strongly typed messages into compact binary streams. Instead of transmitting field names, Protobuf transmits **field tags**.

Each field on the wire begins with a tag byte:
$$\text{Key} = (\text{field\_number} \ll 3) \mid \text{wire\_type}$$

\begin{table}[htbp]
\centering
\small\sffamily
\begin{tabular}{@{}lll@{}}
\toprule
\textbf{Wire Type} & \textbf{Meaning} & \textbf{Used For} \\ \midrule
0 & Varint & \texttt{int32}, \texttt{int64}, \texttt{uint32}, \texttt{bool}, \texttt{enum} \\
1 & 64-bit & \texttt{fixed64}, \texttt{sfixed64}, \texttt{double} \\
2 & Length-delimited & \texttt{string}, \texttt{bytes}, embedded messages, packed arrays \\
5 & 32-bit & \texttt{fixed32}, \texttt{sfixed32}, \texttt{float} \\ \bottomrule
\end{tabular}
\caption{Protocol Buffer wire types.}
\label{tab:protobuf-wire}
\end{table}

\subsubsection{Varints: Variable-Length Quantity Encoding}
Integers in Protobuf are encoded as **Varints**. Each byte reserves its Most Significant Bit (MSB, bit 7) as a **continuation bit**:
\begin{itemize}[leftmargin=15pt]
  \item If $\text{MSB} = 1$, more bytes follow.
  \item If $\text{MSB} = 0$, this is the final byte of the integer.
  \item The remaining 7 bits store the integer in little-endian order.
\end{itemize}
Small numbers (such as $1$) consume **exactly 1 byte** on the wire instead of 4 or 8 bytes in standard C or JSON!

\subsubsection{ZigZag Encoding for Signed Integers}
In two's complement representation, negative numbers like $-1$ have their sign bit set to 1, causing a 64-bit integer to have all 64 bits set, which would force Varint to consume 10 bytes!

Protobuf solves this for \texttt{sint32} and \texttt{sint64} using **ZigZag encoding**, which interleaves positive and negative integers:
$$\text{ZigZag32}(n) = (n \ll 1) \oplus (n \gg 31)$$
This maps $0 \rightarrow 0$, $-1 \rightarrow 1$, $1 \rightarrow 2$, $-2 \rightarrow 3$, ensuring small negative numbers consume only 1 or 2 bytes on the wire!

\subsection{gRPC over HTTP/2 Framing}

gRPC maps procedure calls directly onto HTTP/2 streams:
\begin{itemize}[leftmargin=15pt]
  \item \textbf{Method Path:} Emitted as an HTTP/2 \texttt{POST /<package>.<Service>/<Method>}.
  \item \textbf{Headers:} Includes \texttt{content-type: application/grpc} and \texttt{te: trailers}.
  \item \textbf{Length-Prefixed Message Framing:} Every message payload is preceded by a 5-byte header:
  \begin{itemize}
    \item 1 byte: Compression flag (0 = uncompressed, 1 = compressed with gzip/snappy).
    \item 4 bytes: 32-bit big-endian integer indicating payload length in bytes.
  \end{itemize}
  \item \textbf{Trailers:} The call status is returned in HTTP/2 trailers at the end of the stream: \texttt{grpc-status: 0} (OK), and optional \texttt{grpc-message: "..."}.
\end{itemize}

\begin{lstlisting}[language=C,caption={Protobuf IDL Definition (\texttt{order.proto})}]
syntax = "proto3";

package commerce;

service OrderService {
    // Unary RPC
    rpc CreateOrder (OrderRequest) returns (OrderResponse);
    // Server Streaming RPC
    rpc StreamOrderUpdates (OrderQuery) returns (stream OrderEvent);
}

message OrderRequest {
    int64 customer_id = 1;
    string sku = 2;
    int32 quantity = 3;
}

message OrderResponse {
    string order_id = 1;
    bool success = 2;
}

message OrderQuery {
    string order_id = 1;
}

message OrderEvent {
    string order_id = 1;
    string status = 2;
    int64 timestamp = 3;
}
\end{lstlisting}

% ==========================================
\section{Chapter Summary \& Key Takeaways}

\begin{itemize}[leftmargin=15pt]
  \item \textbf{HTTP/1.1 vs HTTP/2 vs HTTP/3:} HTTP/1.1 suffers from application HoL blocking. HTTP/2 eliminates application HoL blocking via binary frames and stream multiplexing, compressing headers via HPACK, but remains vulnerable to TCP packet-loss HoL blocking. HTTP/3 replaces TCP with QUIC over UDP, delivering independent stream processing and 0-RTT handshakes.
  \item \textbf{REST and HATEOAS:} True REST relies on six architectural constraints. The Richardson Maturity Model traces progression from RPC (Level 0) to hypermedia-driven state engines (Level 3).
  \item \textbf{HTTP Safety and Idempotence:} \texttt{GET} is safe and idempotent; \texttt{PUT} and \texttt{DELETE} are idempotent; \texttt{POST} is non-idempotent.
  \item \textbf{WebSocket Mechanics (RFC 6455):} Upgrades from HTTP/1.1 via a 101 handshake. Mandates 4-byte XOR client masking to prevent proxy cache poisoning attacks.
  \item \textbf{Server-Sent Events (SSE):} Light-weight unidirectional streaming over \texttt{text/event-stream}. Automatic client reconnects leverage \texttt{Last-Event-ID} for deterministic resume.
  \item \textbf{Protobuf Encoding:} Replaces bloated text JSON with binary Varints, ZigZag encoding, and field tag numbers, slashing network serialization overhead in microservice RPC fleets.
\end{itemize}

% ==========================================
\section{Exercises and Real-World Scenarios}

\subsection{Mechanics Exercises}
\begin{enumerate}
  \item \textbf{Exercise 3.1:} Manually encode the integer value $300$ into a Protocol Buffer Varint byte sequence. Show the binary representation, continuation bits, and final hexadecimal bytes.
  \item \textbf{Exercise 3.2:} Explain the exact sequence of HTTP/2 frames exchanged when a client aborts an in-flight file download without closing the underlying TCP connection.
\end{enumerate}

\subsection{Debugging Exercises}
\begin{enumerate}
  \item \textbf{Exercise 3.3:} A microservice fleet migrating from HTTP/1.1 to HTTP/2 behind a corporate NAT firewall reports that during peak morning hours, p99 latency spikes dramatically from 4ms to over 6,000ms whenever network packet loss hits 1.5\%. The operations team suspects CPU saturation on the API servers, but CPU utilization is below 25\%. Diagnose the root architectural cause and provide the remedy.
\end{enumerate}

\subsection{Architecture Design Problem}
\begin{enumerate}
  \item \textbf{Exercise 3.4:} Architect a high-frequency cryptocurrency order book feed that serves 500,000 concurrent browser clients. Detail the transport protocol choice (WebSocket vs SSE vs HTTP/2 polling), horizontal connection scaling layer, and backpressure strategy when slow clients cannot consume messages at wire speed.
\end{enumerate}

% ==========================================
\section{Detailed Solutions}

\subsection{Solution 3.1: Protobuf Varint Encoding of 300}
\begin{enumerate}[leftmargin=15pt]
  \item Binary representation of $300$: $100101100_2$ (9 bits).
  \item Split into 7-bit chunks (least significant first):
  \begin{itemize}
    \item Lower 7 bits: $0101100_2$ (decimal 44).
    \item Upper bits: $0000010_2$ (decimal 2).
  \end{itemize}
  \item Set continuation bit (MSB):
  \begin{itemize}
    \item Byte 1 has more data to follow $\rightarrow \text{MSB} = 1$: $10101100_2 = \texttt{0xAC}$.
    \item Byte 2 is the last byte $\rightarrow \text{MSB} = 0$: $00000010_2 = \texttt{0x02}$.
  \end{itemize}
  \item Final encoded Varint bytes: \texttt{0xAC 0x02}.
\end{enumerate}

\subsection{Solution 3.2: HTTP/2 Stream Cancellation}
The client emits a \texttt{RST\_STREAM} frame (type \texttt{0x03}) specifying the target stream ID and error code \texttt{CANCEL (0x08)}. The server receives the frame, halts processing on that specific stream, frees local buffers, and continues processing all other active multiplexed streams across the shared TCP socket without disturbance.

\subsection{Solution 3.3: Diagnosing TCP Head-of-Line Blocking}
\textbf{Root Cause:} The fleet is suffering from **TCP-level Head-of-Line (HoL) blocking** in HTTP/2. Because HTTP/2 multiplexes hundreds of logical streams onto a single TCP connection, a 1.5\% packet loss causes the Linux kernel to hold back all received packets in the TCP queue while waiting for retransmission of the lost segment. Even though CPU utilization is low, all streams sharing the TCP connection are completely frozen.

\textbf{Remedy:}
\begin{enumerate}
  \item Migrate the transport layer to **HTTP/3 (QUIC)**, which operates over UDP and guarantees independent stream processing.
  \item In the interim, tune the client/proxy connection pool to maintain a small pool of 4--6 HTTP/2 connections per origin instead of multiplexing everything onto a single TCP socket.
\end{enumerate}

\subsection{Solution 3.4: Order Book Feed Architecture}
See Section 3.4 and Section 3.5. Because order book updates flow exclusively from the trading engine to clients, **Server-Sent Events (SSE)** over HTTP/2 (or HTTP/3) combined with a Redis Pub/Sub backplane provides the most resilient architecture. Clients connect to stateless Edge Gateways (e.g., Envoy or Nginx). For slow clients, the gateway implements a ring buffer per connection that drops stale intermediate quotes while preserving sequence snapshots.

% ==========================================
\section{Further Reading \& Citations}

\begin{itemize}[leftmargin=15pt]
  \item Fielding, R. T. (2000). \textit{Architectural Styles and the Design of Network-based Software Architectures}. Doctoral dissertation, University of California, Irvine.
  \item Belshe, M., Peon, R., \& Thomson, M. (2015). \textit{Hypertext Transfer Protocol Version 2 (HTTP/2)}. RFC 7540.
  \item Iyengar, J., \& Thomson, M. (2021). \textit{QUIC: A UDP-Based Multiplexed and Secure Transport}. RFC 9000.
  \item Fette, I., \& Melnikov, A. (2011). \textit{The WebSocket Protocol}. RFC 6455.
  \item Google. (2023). \textit{Protocol Buffers Developer Guide \& Encoding Internals}. \url{https://protobuf.dev/}
\end{itemize}
